\clearpage\pagestyle{otMain}\setlength{\otBlocEspacement}{2.5mm}\fontsize{11}{15}\selectfont\corrigesChapitre{1}{Opérations et priorités}
\begin{blocLivre}{cor-1-1}\begin{corrige}{1}{Addition décimale}\demarche{Aligner les unités.}
Aligner les rangs : $4{,}70+2{,}85=\res{7{,}55}$. Le total est proche de $5+3=8$.
\end{corrige}\end{blocLivre}
\begin{blocLivre}{cor-1-2}\begin{corrige}{2}{Différence décimale}\demarche{Vérifier par addition.}
$10{,}00-3{,}48=\res{6{,}52}$. Contrôle : $6{,}52+3{,}48=10$.
\end{corrige}\end{blocLivre}
\begin{blocLivre}{cor-1-3}\begin{corrige}{3}{Produit de décimaux}\demarche{Utiliser les rangs décimaux.}
$8\times15=120$. Les deux facteurs ont ensemble deux chiffres après la virgule : le produit est $\res{1{,}20=1{,}2}$.
\end{corrige}\end{blocLivre}
\begin{blocLivre}{cor-1-4}\begin{corrige}{4}{Diviser par un décimal}\demarche{Rendre le diviseur entier.}
On multiplie les deux nombres par 10 : $96:4=\res{24}$. Contrôle : $24\times0{,}4=9{,}6$.
\end{corrige}\end{blocLivre}
\begin{blocLivre}{cor-1-5}\begin{corrige}{5}{Nommer une expression}\demarche{Repérer la structure avant le calcul.}
Le premier calcul est une somme : $5+14=\res{19}$. Le second est un produit : $7\times7=\res{49}$.
\end{corrige}\end{blocLivre}
\begin{blocLivre}{cor-1-6}\begin{corrige}{6}{Même priorité}\demarche{Garder l’ordre à priorité égale.}
De gauche à droite : $36:6\times3=6\times3=\res{18}$ et $24-8+5=16+5=\res{21}$.
\end{corrige}\end{blocLivre}
\begin{blocLivre}{cor-1-7}\begin{corrige}{7}{Deux étages de parenthèses}\demarche{Commencer à l’intérieur.}
La parenthèse donne 3 ; le crochet vaut $3\times3+7=16$. Donc $42-16=\res{26}$.
\end{corrige}\end{blocLivre}
\begin{blocLivre}{cor-1-8}\begin{corrige}{8}{Distribuer pour gagner du temps}\demarche{Chercher une somme de facteurs simple.}
$23(100+1)=2300+23=\res{2323}$. On regroupe le second calcul : $48(9{,}5+0{,}5)=48\times10=\res{480}$.
\end{corrige}\end{blocLivre}
\begin{blocLivre}{cor-1-9}\begin{corrige}{9}{Une division à contrôler}\demarche{Multiplier les deux nombres par 100.}
$6{,}3:0{,}15=630:15=\res{42}$. Le contrôle donne $42\times0{,}15=6{,}3$.
\end{corrige}\end{blocLivre}
\begin{blocLivre}{cor-1-10}\begin{corrige}{10}{Corriger une priorité}\demarche{Conserver l’expression de départ.}
Il a additionné avant de multiplier, sans parenthèses. Il faut $7+15=\res{22}$. Le résultat 50 correspond à $(7+3)\times5$.
\end{corrige}\end{blocLivre}
\begin{blocLivre}{cor-1-11}\begin{corrige}{11}{Une sortie}\demarche{Déterminer d’abord ce que paie une personne.}
Chaque élève paie $6{,}50+3{,}20=9{,}70$ euros. Le total est $4(6{,}50+3{,}20)=\res{38{,}80}$ euros. Distribuer donne aussi $26+12{,}80$.
\end{corrige}\end{blocLivre}
\begin{blocLivre}{cor-1-12}\begin{corrige}{12}{Des bouteilles}\demarche{Interpréter le quotient.}
$8{,}4:0{,}75=11{,}2$ : on remplit $\res{11}$ bouteilles. Elles contiennent $11\times0{,}75=8{,}25$ L ; il reste $\res{0{,}15\ \mathrm L}$. Une douzième bouteille n’est pas pleine.
\end{corrige}\end{blocLivre}
\begin{blocLivre}{cor-1-13}\begin{corrige}{13}{Un budget partagé}\demarche{Distinguer dépense, reste et partage.}
Les achats coûtent $5\times12{,}40=62$ euros. Le reste est $125-62=63$ euros. Chaque atelier reçoit $(125-5\times12{,}40):6=\res{10{,}50}$ euros.
\end{corrige}\end{blocLivre}
\begin{blocLivre}{cor-1-14}\begin{corrige}{14}{Une expression à inventer}\demarche{Chercher une étape intermédiaire.}
Une solution est $\res{(8+6)\times2=28}$. Sans parenthèses, $8+6\times2=20$. Les parenthèses imposent de calculer la somme avant le produit.
\end{corrige}\end{blocLivre}
\begin{blocLivre}{cor-1-15}\begin{corrige}{15}{Comparer deux offres}\demarche{Traduire chaque offre séparément.}
Premier magasin : $10\times2{,}80=28$ euros. Second : $2\times12+1{,}50=25{,}50$ euros. Le second coûte $\res{2{,}50}$ euros de moins. Les frais sont ajoutés une seule fois.
\end{corrige}\end{blocLivre}
